\documentstyle[%


righttag%


,11pt%


,amssymb%


,russian%               remove in Eng.


,izvru%                 Set izven% in Eng.


]{amsart}





%





\newcommand{\tts}[1]{}





%\hoffset=-15mm%                  For Eng.


\setcounter{page}{39}


\god{2002}


\nomer{1~(476)} %                        Remove in Eng.


%\nomer{Vol.\,46, No.\,1}%               Set in Eng.


%\str{\pageref{begin}--\pageref{end}}%  Set in Eng.


\udk{517.544}





\author{\it Б.А.\,Кац}


% NB:   ^^ remove in Eng.


\title{Об обобщении одной теоремы Н.А.\,Давыдова}


\thanks{Работа выполнена при финансовой поддержке Российского фонда


фундаментальных исследований, грант \No\,01-01-00088.}





\date{}


%\pagestyle{myheadings}%         Set in Eng.


\pagestyle{plain} %              Remove in Eng.


\begin{document}


%\thispagestyle{plain}%          Set in Eng.


\maketitle


\label{begin}








Пусть $\Gamma$ есть простая замкнутая спрямляемая кривая на комплексной


плоскости, разбивающая ее на области $D^+$ и $D^{-}\ni\infty$. Тогда для


любой заданной на $\Gamma$ непрерывной функции $f(t)$ определен интеграл


типа Коши


\begin{equation}


K_{f}(z) = \frac{1}{2\pi i}\int_{\Gamma}{\frac{f(t)dt}{t-z}}.


\end{equation}


Он представляет голоморфную в $\overline{\Bbb C}\setminus\Gamma$ функцию.


Традиционный интерес вызывает вопрос о граничных значениях этой функции.


В значительной степени это обусловлено приложениями интеграла типа Коши при


решении краевых задач и сингулярных интегральных уравнений \cite{Ga},


\cite{Mu}. В простейшем случае это вопрос о возможности непрерывного


продолжения функции $K_{f}(z)$ на кривую $\Gamma$ слева (из области $D^+$) и


справа (из $D^-$).





Хорошо известно, что интеграл (1) по кусочно-гладкой кривой $\Gamma$


непрерывен в замыканиях $D^+$ и $D^-$, если его плотность $f(t)$


удовлетворяет условию Г\"ельдера


\begin{equation}


\sup \bigg\{\frac{\left| f(t')-f(t'')\right|}{\left| t'-t''\right|^{\nu}}


: t',t''\in \Gamma,\ t'\ne t''\bigg\}\equiv h_{\nu}(f,\Gamma)<\infty


\end{equation}


с любым показателем $\nu\in (0,1]$. Как отмечается в


\cite{Mu}, этот факт был известен еще Гарнаку, Морера и Сохоцкому.





Ниже через $H_{\nu}(\Gamma)$ обозначим пространство Г\"ельдера, т.\,е.


множество всех заданных на $\Gamma$ функций, удовлетворяющих условию (2).


При $\nu = 1$ --- это класс Липшица.





Одним из первых вопрос о непрерывности интеграла типа Коши по негладкой


кривой исследовал Н.А.\,Давыдов. В частности, он установил следующий факт.


\medskip





{\bf Теорема Н.А.\,Давыдова} \cite{D}. {\it


Интеграл типа Коши $(1)$ по замкнутой


спрямляемой кривой~$\Gamma$ непрерывен в замыканиях областей $D^+$ и $D^-$,


на которые эта кривая разбивает комплексную плоскость, если его плотность


$f$ принадлежит классу Липшица $H_{1}(\Gamma)$.


}


\medskip





В данной статье теорема Н.А.\,Давыдова обобщается на неспрямляемые кривые.


Прежде всего уточним, в каком смысле понимается интеграл (1),


если кривая $\Gamma$ неспрямляема. Приведем соответствующее


определение.





Зафиксируем положительное направление обхода кривой $\Gamma$ (всюду ниже


это направление против часовой стрелки) и выберем точку $t^*\in\Gamma$,


которая будет служить началом и концом этого обхода. Тем самым


задаем на $\Gamma\setminus t^*$ естественные отношения порядка


$z_1\prec z_2$, если $z_1$ предшествует $z_2$ при таком обходе,


и $z_1\preceq z_2$, если $z_1 \prec z_2$ или $z_1 = z_2$. Множество


$\Gamma\setminus t^*$ можно пополнить наибольшим и наименьшим элементами


$t_{\min}$ и $t_{\max}$; оба эти элемента могут быть в понятном


смысле отождествлены с точкой $t^*$.





Последовательность точек $\{z_j\}_{j=0}^{j=n}$ назовем разбиением


$\Gamma$, если $t_{\min}=z_0\prec z_1\prec z_2\prec\dotsb\prec z_n=t_{\max}$,


имея в виду, что эти точки разбивают $\Gamma$ на дуги $\gamma_j$ так, что


началом $\gamma_j$ является точка $z_{j-1}$, а концом --- точка $z_j$,


$j = 1,2,\dotsc, n$. Если кроме этого заданы точки $\{w_j\}_{j=1}^{j=n}$


такие, что $z_{j-1}\preceq w_j\preceq z_j$, $j=1,2,\dotsc,n$, то разбиение


называется пунктированным.





\begin{defn}


Пусть $\Gamma$ есть простая замкнутая кривая с фиксированными направлением


обхода и ``началом--концом'' $t^*$, на которой заданы функции $f(z)$ и $g(z)$.


Рассмотрим всевозможные пунктированные разбиения $\Gamma$, и с каждым из них


свяжем сумму Римана--Стильтьеса


$S = \sum\limits_{j=1}^{n} f(w_j)(g(z_j) - g(z_{j-1}))$.


Если для любого $\varepsilon > 0$ существует такое $\delta >0$, что


$\left|S-J\right| <\varepsilon$ для любого пунктированного разбиения,


удовлетворяющего условию


$\left|z_j - z_{j-1}\right| < \delta$, $j=1,2,\dots,n$,


то число $J$ есть контурный интеграл Стильтьеса: $J=\int\limits_{\Gamma}fdg$.


\end{defn}





Нетрудно убедиться, что ни существование этого интеграла, ни


его значение не зависят от выбора точки $t^*$.





Существует простая связь этого интеграла с обычным интегралом Стильтьеса по


отрезку вещественной оси. Пусть $z=z(x)$ --- любое непрерывное строго


монотонное (в смысле отношения~$\prec$) отображение отрезка $I = [0, 1]$


на кривую $\Gamma$ такое, что $z(0)=t_{\min}$ и $z(1)=t_{\max}$.


Очевидно,


$\int\limits_{\Gamma}fdg=\int\limits_0^1 f(z(x))dg(z(x))$, где в правой части


стоит обычный интеграл Стильтьеса по отрезку~$I$. Конечно, отображение


$z(x)$ не единственно, но величина интеграла и его существование не зависят


от выбора этого отображения.





Теперь обратимся к условиям существования интеграла Стильтьеса в


случаe $g(z)\equiv z$, $g(z(x))\equiv z(x)$.


Hаиболее известное условие, предполагающее ограниченность


вариации функции~$g$ (напр., \cite{R}), не может


быть использовано, поскольку


ограниченность вариации отображения $z(x)$ равносильна спрямляемости его


образа $\Gamma$. Вместо него воспользуемся результатом


\cite{Y2}, основанным на следующем обобщении понятия вариации.





\begin{defn}


Пусть $\Phi (x)$ есть заданная при $x\ge 0$ непрерывная


возрастающая функция такая, что $\Phi (0)=0$. Класс $V_{\Phi}(\Gamma)$


состоит из всех заданных на кривой $\Gamma$ функций $f(z)$, удовлетворяющих


условию


$$


\sup_{\tau}\sum_{j=1}^{n}\Phi(\left|f(z_{j})-f(z_{j-1})\right|)


\equiv v_{\Phi}(f;\Gamma)<\infty;


$$


здесь точная верхняя грань берется по


всем разбиениям $\tau = \{z_{j}\}_{j=0}^{j=n}$.


Если $\Phi (x) = x^p$, $p\ge 1,$ то этот класс обозначается через


$V_{p}(\Gamma)$.


\end{defn}





Обычно эти классы рассматривают на отрезке вещественной оси $I$,


используя при этом обычные неравенства вместо $\prec$ и $\preceq$ и


полагая в определении разбиений $t_{\min} = 0$, $t_{\max} = 1$. Очевидно,


$V_{1}(I)$ есть обычный класс функций ограниченной вариации. Класс


$V_{2}(I)$ был введен Н.\,Винером \cite{W}. В общем виде классы $V_{p}(I)$ и


$V_{\Phi}(I)$ были введены и использованы при изучении интеграла


Стильтьеса в \cite{Y1} и \cite{Y2} соответственно.


\medskip





{\bf Теорема Л.\,Юнга} \cite{Y2}. {\it


Интеграл Стильтьеса


$\int\limits_0^1{fdg}$ существует при условиях $f\in V_{\Phi}(I)$,


$g\in V_{\Psi}(I)$, если функции $f$ и $g$ не имеют общих точек разрыва


и сходится ряд Юнга$:$


$$


\sum_{n=1}^{\infty}\varphi (1/n)\psi (1/n) <\infty,


$$


где $\varphi$ и $\psi$ --- функции, обратные к $\Phi$ и $\Psi$ соответственно.


}


\medskip





Чтобы применять эту теорему для обоснования существования контурного


интеграла Стильтьеса $\int\limits_{\Gamma}{fdz}$, придадим условию


$z(x)\in V_{\Phi}(I)$ геометрическую форму.





\begin{defn}


Будем относить кривую $\Gamma$ к классу $R_{\Phi}$ и называть ее


$\Phi$-спрямляемой, если величина


$$


\sigma_{\Phi}(\Gamma) \equiv


\sup_{\tau}\sum_{j=1}^{n}\Phi(\left|z_{j} - z_{j-1}\right|)


$$


конечна; точная верхняя грань здесь берется по всем


разбиениям $\tau = \{z_{j}\}_{j=0}^{j=n}$.


Если $\Phi (x) = x^p$, $p\ge 1,$ то будем говорить о $p$-спрямляемости


и писать $R_p$ и $\sigma_p$ вместо $R_{\Phi}$ и $\sigma_{\Phi}$


соответственно.


\end{defn}





При $p>1$ класс $R_p$ содержит неспрямляемые кривые. Например,


известная неспрямляемая ``снежинка'' фон Коха (напр., \cite{F})


является $(\log_{3}{4})$-спрямляемой.





Теперь можно сформулировать обобщение теоремы Н.А.\,Давыдова для


неспрямляемых кривых.





\begin{thm}\label{phi}


Если кривая $\Gamma$ является $\Phi$-спрямляемой,


где функция $\Phi (x)$


выпуклая\,\footnote{т.\,е. $\Phi (\alpha x + \beta y)\le \alpha \Phi (x) +


\beta \Phi (y)$ при $x,y>0$, $\alpha,\beta\ge0$, $\alpha+\beta = 1$.},


и $f\in H_{1}(\Gamma)$, то интеграл типа 


Коши $(1)$ существует и непрерывен в замыканиях областей $D^+$


и $D^-$ при условии


\begin{equation}


\sum_{n=1}^{\infty}{\varphi^{2}(1/n)} < \infty,


\end{equation}


где $\varphi$ --- функция, обратная к $\Phi$.


\end{thm}





{\bf Доказательство} начнем с обоснования существования контурного интеграла


(1) при условиях теоремы \ref{phi}. Как уже отмечалось, $\Phi$-спрямляемость


кривой $\Gamma$ равносильна включению отображения $z=z(x):I\mapsto\Gamma$ в


класс $V_{\Phi}(I)$. Пусть функция $f\in H_{1}(\Gamma)$ отличнa от постоянной,


т.\,е. $h_{1}(f;\Gamma) = h\not= 0.$ Положим $f_{0}(t)\equiv f(t)/h$. Тогда


для любого разбиения $\{z_j\}_{j=0}^{j=n}$ имеем


$$


\sum_{j=1}^{n}\Phi(\left|f_{0}(z_{j})-f_{0}(z_{j-1})\right|)\le


\sum_{j=1}^{n}\Phi(|z_{j}-z_{j-1}|) \le \sigma_{\Phi}(\Gamma),


$$


т.\,е. $f_{0}\in V_{\Phi}(\Gamma)$, $f_{0}(z(x))\in V_{\Phi}(I)$,


и $v_{\Phi}(f_{0}\circ z; I)\le \sigma_{\Phi}(\Gamma)$.


Тогда в силу теоремы Л.\,Юнга интеграл $\int\limits_0^1{f_{0}(z(x))dz(x)} =


\int\limits_{\Gamma}{f_{0}dz}$ существует при условии (3). Очевидно, тогда


существует и интеграл


$\int\limits_{\Gamma}{fdz}$. При любом фиксированном $z\in


{\Bbb C}\setminus\Gamma$ функция $\wt f(t) = f(t)/(t - z)$ удовлетворяет


условию Липшица вместе с $f(t)$. Поэтому интеграл типа Коши (1) также


существует при условии (3).





Перейдем к доказательству непрерывности этого интеграла на $\Gamma$.


Продолжим функцию $f$ с кривой $\Gamma$ на всю плоскость ${\Bbb C}$


посредством оператора продолжения Уитни ${\cal E}_0$ \cite{St} и


обозначим полученную функцию через $f^w(z)$. Напомним некоторые свойства


продолжения Уитни липшицевой функции, доказательства которых можно найти в


\cite{St},


\begin{itemize}


\item[а)] сужение $f^w$ на $\Gamma$ совпадает с $f$;


\item[б)] $f^{w}\in H_{1}({\Bbb C})$;


\item[в)] в ${\Bbb C}\setminus\Gamma$ продолженная функция $f^w$ имеет частные


производные всех порядков;


\item[г)] первые частные производные функции $f^w$ ограничены.


\end{itemize}





Докажем, что для функции $f^w$ справедлива формула Стокса


\begin{equation}


\int\limits_{\Gamma}f^{w}(z)dz = -\iint\limits_{D^+}


\frac{\partial f^w}{\partial\overline z}dz\,d\overline z.


\end{equation}


Для этого фиксируем отображение $z=z(x):I\mapsto\Gamma$, о котором шла речь


выше. Положим


$x_{j,n} = j/2^{n}$, $z_{j,n} = z(x_{j,n})$, $j = 0,1,\dotsc,2^n$,


$n = 1,2,\dotsc$, и определим комплекснозначную кусочно-линейную функцию


$z_{n}(x)$ равенствами $z_{n}(x_{j,n}) = z(x_{j,n})$, $j=0,1,\dotsc,2^n$;


на интервалах $\big[\frac{j}{2^{n}},\frac{j+1}{2^n}\big]$ эта функция линейна.


Она отображает $I$ на замкнутую ломаную (цепь) $\Gamma_n$, возможно,


самопересекающуюся, которая ограничивает один или несколько многоугольников.


Совокупность этих многоугольников с приписанным каждому из них


знаком $+$ или $-$


в зависимости от того, слева или справа от $\Gamma_n$ этот многоугольник


находится, будем обозначать $D^+_n$. Далее, положим $f^{w}_{0}(z)\equiv


f^{w}(z)/h_{1}(f^w; {\Bbb C})$. Очевидно, для этой функции справедлива


формула Стокса


\begin{equation}


\int\limits_{\Gamma_n}f^{w}_{0}(z)dz = -\iint\limits_{D^+_n}


\frac{\partial f^w_0}{\partial\overline z}dz\,d\overline z,


\end{equation}


если в правой части стоит сумма интегралов во вышеописанным многоугольникам


с соответствующими знаками. Левую часть этого равенства можно записать в


виде интеграла Стильтьеса


$\int\limits_{0}^{1}{f^{w}_{0}(z_{n}(x))dz_{n}(x)}$.


В силу выпуклости $\Phi$ имеем $V_{\Phi}(z_{n};I)\le V_{\Phi}(z;I) =


\sigma_{\Phi}(\Gamma)$ (см. \cite{LO}; до этого условие выпуклости $\Phi$


в данном доказательстве не использовалось). Так же, как и выше, отсюда


следует оценка


$v_{\Phi}(f^{w}_{0}\circ z_n; I)\le \sigma_{\Phi}(\Gamma)$. Таким образом,


последовательности функций $\{z_{n}(x)\}$ и $\{f^{w}_{0}(z_{n}(x))\}$


равномерно сходятся и имеют ограниченные в совокупности $\Phi$-вариации.


Поэтому (см. \cite{LO}; здесь также используется выпуклость функции $\Phi$)


предел последовательности интегралов


$\int\limits_{0}^{1}{f^{w}_{0}(z_{n}(x))dz_{n}(x)}$ существует и равен


$\int\limits_{0}^{1}{f^{w}_{0}(z(x))dz(x)}=\int\limits_{\Gamma}{f^w_0 dz}$.


Легко убедится, что предел интеграла в правой части (5) также существует и


равен


$\iint\limits_{D^+}{\frac{\partial f^w_0}{\partial\ov z}dz\,d\ov z}$.


Отсюда следует справедливость равенства (4).





Из только что проведенного рассуждения следует, что формула Стокса (4)


справедлива не только для продолжения Уитни, но и для любой функции $f^w$,


удовлетворяющая условиям а)--г).


Более того, эта функция может быть определена не в ${\Bbb C}$, а в


любой содержащей $D^+$ области $D'$; соответственно условие б) можно


заменить на $f^{w}\in H_{1}(D')$, а в условииях в) и г) можно требовать


существования и ограниченности


первых частных производных лишь в $D'\setminus\Gamma$.





Применим теперь эту формулу к функции $f(t)/(t - z)$. Если фиксированная


точка $z$ лежит в $D^-$, то подходящим ее продолжением является функция


$f^{w}(t)/(t - z)$, $t\in D'$, где $f^w$ --- продолжение Уитни; в качестве


области $D'$ здесь можно взять любую конечную область, содержащую $D^+$ и не


содержащую $z$. В результате получим


$$


\int\limits_{\Gamma}{\frac{f(t)dt}{t-z}} = -


\iint\limits_{D^{+}}


{\frac{\partial f^{w}}{\partial\overline{t}}\frac{dt\,d\overline{t}}{t-z}},


\quad z\in D^-.


$$


Если $z\in D^-$, то продолжение $f^{w}(t)/(t - z)$ следует заново


определить в малом круге $\{t:|t-z|<\varepsilon\}$ так, чтобы оно стало


непрерывно дифференцируемым в $D^+$. Рутинные вычисления дают


в итоге


$$


\int\limits_{\Gamma}{\frac{f(t)dt}{t-z}} = 2\pi i f^{w}(z) -


\iint\limits_{D^{+}}


{\frac{\partial f^{w}}{\partial\overline{t}}\frac{dt\,d\overline{t}}{t-z}},


\quad z\in D^+.


$$


Объединив две последние формулы, получаем


\begin{equation}


K_{f}(z) = \chi (z)f^{w}(z) - \frac{1}{2\pi i}


\iint\limits_{D^{+}}


{\frac{\partial f^{w}}{\partial\overline{t}}\frac{dt\,d\overline{t}}{t-z}},


\quad z\in{\Bbb C}\setminus\Gamma,


\end{equation}


где $\chi (z)$ --- характеристическая функция области $D^+$.


Интеграл в правой части (6) подробно изучен (см., напр., \cite{V}).


Если его плотность ограничена, то он представляет функцию, непрерывную во


всей плоскости. Непрерывность слагаемого $\chi (z)f^{w}(z)$ в замыканиях


областей $D^+$ и $D^-$ очевидна.$\quad\square$


\medskip





{\bf Следствие.}


Если кривая $\Gamma$ является $p$-спрямляемой, $1\le p<2$, и


$f\in H_{1}(\Gamma)$, то интеграл типа Коши (1) существует и непрерывен в


замыканиях областей $D^+$ и $D^-$.


\medskip





Функция $\Phi (x) = x^p$, $p\ge 1$, выпукла, и условие (3) для нее


имеет вид $p<2.$ Таким образом, следствие вытекает непосредственно


из теоремы \ref{phi}.





Pассмотрим более простое условие существования интеграла


Стильтьеса $\int\limits_0^1{fdg}$, чем теорема Юнга. Хорошо известно (см.,


напр., \cite{R}), что этот интеграл существует, если функция


$g$ непрерывна, а $f\in V_{1}(I)$. Следовательно, интеграл


$\int\limits_{\Gamma}{fdz}$ существует, если $\Gamma$ есть произвольная


жорданова кривая, а $f\in V_{1}(\Gamma)$. Если кривая $\Gamma$ неспрямляема,


то класс Липшица $H_{1}(\Gamma)$ не является подмножеством класса


$V_{1}(\Gamma)$. В связи с этим возникает следующее предположение:





-- если $f\in H_{1}(\Gamma)\cap V_{1}(\Gamma)$, то интеграл типа Коши (1) по


произвольной замкнутой жордановой кривой $\Gamma$ непрерывен в


замыканиях областей $D^+$ и $D^-$.





Автору пока не удалось доказать это предположение в полном объеме, но один


результат в этом направлении здесь приведем.





\begin{defn}


Класс ${\cal G\cal C}$ состоит из простых замкнутых жордановых кривых,


представимых в виде объединения конечного числа графиков (возможно,


повернутых) непрерывных вещественных функций.


\end{defn}





Пусть $\Gamma\in {\cal G\cal C}$, т.\,е.


эта кривая представима в виде объединения


графиков $\gamma_j$, $j=1,2,\dotsc,m$, $m=m(\Gamma)$,


непрерывных вещественных функций $y=y_{j}(x)$, 


каждый из которых повернут на угол $\theta_j$, $j=1,2,\dotsc,m$.


Тогда $\gamma_j = \{t=(x+iy)e^{i\theta_j}:a_j < x < b_j,\ y=y_{j}(x)\}$,


где $[a_j, b_j]\equiv I_j$ --- отрезок вещественной оси,


и сужение на $\gamma_j$ любой заданной на $\Gamma$ функции $f(t)$ можно


рассматривать как функцию вещественной переменной $x$:


$f\vert_{\gamma_j}(e^{i\theta_j}(x+iy_{j}(x)))\equiv f_{j}(x)$,


$x\in I_j$.





\begin{defn}


Заданную на кривой $\Gamma\in {\cal G\cal C}$ функцию $f(t)$ будем относить


к классу ${\cal H}_{1}(\Gamma)$, если она непрерывна и


$f_{j}\in H_{1}(I_j)$ при $j=1,2,\dotsc, m(\Gamma)$.


\end{defn}





Нетрудно убедиться, что ${\cal H}_{1}(\Gamma)\subset


H_{1}(\Gamma)\cap V_{1}(\Gamma)$, причем на неспрямляемой


кривой $\Gamma$ это включение строгое.





\begin{thm}\label{arbitr}


Если $\Gamma\in{\cal G\cal C}$ и $f\in {\cal H}_{1}(\Gamma)$, то


интеграл типа Коши $(1)$ непрерывен в замыканиях областей $D^+$ и $D^-$.


\end{thm}





\begin{pf}


Существование контурного интеграла $\int\limits_{\Gamma}{fdz}$


и интеграла типа Коши при условиях этой теоремы очевидно. Чтобы получить


аналог формулы (6), рассмотрим сначала область вида


$D = \{z=x+iy: a<x<b,\ 0<y<y(x)\}$, где $y(x)$ --- положительная непрерывная


функция. Любая заданная на отрезке $J=[a,b]$ функция $f\in H_{1}(J)$


допускает тривиальное продолжение в полосу $J^* = J\times{\Bbb R}\supset D$


до функции $f^{*}(x+iy) = f(x)$. Очевидно, $f^*\in W^{1}_{\infty}$, и для


этой функции справедлива формула Стокса в любой области $D'\subset J^*$


с кусочно-гладкой границей. Пусть $\{y_{n}(x)\}$ есть последовательность


положительных гладких функций, равномерно сходящаяся к $y(x)$ на отрезке


$J$. Положим $D_n = \{z=x+iy: a<x<b,\ 0<y<y_{n}(x)\}$. Тогда


$$


\int\limits_{\partial D_n}{f^{*}(z)dz} = -\iint\limits_{D_n}


\frac{\partial f^*}{\partial\overline z}dz\,d\overline z.


$$


Последовательность функций $\{f^{*}(y_{n}(x)\}$ равномерно сходится к


$f^{*}(y(x))$, имея ограниченные в совокупности вариации. В силу стандартных


теорем о переходе к пределу под знаком интеграла Стильтьеса (см., напр.,


\cite{R}) отсюда следует справедливость формулы Стокса


$$


\int\limits_{\partial D}f^{*}(z)dz = -\iint\limits_{D}


\frac{\partial f^*}{\partial\overline z}dz\,d\overline z.


$$


Как и при доказательстве теоремы \ref{phi}, из этой формулы получаем


\begin{equation}


\frac{1}{2\pi i}\int\limits_{\partial D}{\frac{f^{*}(t)dt}{t-z}} =


\chi (z)f^{*}(z) - \frac{1}{2\pi i}


\iint\limits_{D}


{\frac{\partial f^{*}}{\partial\overline{t}}\frac{dt\,d\overline{t}}{t-z}},


\quad z\in{\Bbb C}\setminus\partial D,


\end{equation}


где $\chi (z)$ --- характеристическая функция области $D$, а из этого


равенства следует непрерывность интеграла типа Коши, стоящего в его левой


части, в замыканиях области $D$ и ее дополнения. Складывая равенства


вида (7) для областей, связанных с составляющими кривую $\Gamma$ графиками


$\gamma_j$, после простых преобразований получаем утверждение теоремы.


\end{pf}





Отметим, что из доказательств теорем \ref{phi},~\ref{arbitr} следует, что


при их условиях интеграл типа Коши (1) обладает и некоторыми дополнительными


свойствами. Так, из формул (6) и (7) следует, что разность его граничных


значений на $\Gamma$ равна $f(t)$, что позволяет использовать этот интеграл


при решении краевых задач подобно тому, как это делается в случае спрямляемых


кривых (см. \cite{Ga},~\cite{Mu}). Далее, из оценок \cite{V} следует, что эти


граничные значения принадлежат классу $H_{1-\varepsilon}(\Gamma)$ для сколь


угодно малого $\varepsilon > 0$. Этот факт также может оказаться полезным


при исследовании краевых задач в областях с неспрямляемыми (в


том числе с фрактальными) границами.








\begin{thebibliography}{11}





\bibitem{Ga}


Гахов Ф.Д. {\it Краевые задачи.} -- М.: Наука, 1977. -- 640 с.


\bibitem{Mu}


Мусхелишвили Н.И. {\it Сингулярные интегральные уравнения.


Граничные задачи теории функций и некоторые их приложения


к математической физике}. -- М.: Физматгиз, 1962. -- 600 с.


\bibitem{D}


Давыдов Н.А. {\it Некоторые вопросы теории граничных значений аналитических


функций}: Дис. $\dotsc$ канд. физ.-матем. нaук. -- Москва, 1949.


\bibitem{R}


Рудин У. {\it Основы математического анализа.} -- М.: Мир, 1976. -- 320 с.


\bibitem{Y2}


Young L.C. {\it General inequalities for Stieltjes integrals and the


convergence of Fourier series} // Math. Annalen, Berlin. -- 1938. --


Bd.\,115. -- H.\,4. -- S.\,581--612.


\bibitem{W}


Wiener N. {\it The quadratic variation of a function and its Fourier


coefficients} // Massachusetts J. of Math. -- 1924. -- V.\,3. -- P.\,72--94.


\bibitem{Y1}


Young L.C. {\it An inequality of the H\"older type, connected with Stieltjes


integration} // Acta Math., Uppsala. -- 1936. --


Bd.\,36. -- \No\,3--4. -- S.\,251--282.


\bibitem{F}


Федер Е. {\it Фракталы.} -- М.: Мир, 1991. -- 280 с.


\bibitem{St}


Стейн И. {\it Сингулярные интегралы и дифференциальные свойства


функций.} -- М.: Мир, 1973. -- 342 с.


\bibitem{LO}


Lesniewicz R., Orlicz W. {\it On generalized variation} II // Studia


Math. -- 1973. -- V.\,45. -- \No\,1. -- P.\,71--109.


\bibitem{V}


Векуа И.Н. {\it Обобщенные аналитические функции.}


-- М.: Наука, 1988. -- 512 с.





\end{thebibliography}





\vskip 2cm





\post{\it Казанская государственная}{\it Поступила}


\post{\it архитектурно-строительная академия}{12.04.1999}





\label{end}


\end{document}


